\section{Sylow 子群的存在性}


\subsection{先固定记号}


设 $p$ 是素数，有限群 $G$ 的阶写成

\[
|G|=p^a m,\qquad p\nmid m.
\]

若 $a\ge 1$，阶为 $p^a$ 的子群称为 \textbf{Sylow $p$-子群}；记所有 Sylow $p$-子群组成的集合为
$\operatorname{Syl}_p(G)$，其个数记作 $n_p$。这里的 $p^a$ 是 $|G|$ 中能出现的最大 $p$-幂，因此 Sylow 子群是“尽可能大的 $p$-子群”，而不是任意一个阶为 $p$ 的子群。

当 $a=0$ 时，$p\nmid |G|$，只有平凡子群 $\{e\}$，把它视为唯一的 Sylow $p$-子群即可。下面只讨论 $a\ge1$。

\subsection{定理：Sylow 子群存在}


\textbf{Sylow 第一定理（存在性部分）}：若 $|G|=p^a m$ 且 $p\nmid m$，则 $G$ 至少有一个阶为 $p^a$ 的子群。

证明采用 $|G|$ 的归纳。需要用到第 3 节的类方程

\[
|G|=|Z(G)|+\sum_i [G:C_G(x_i)],
\]

其中 $x_i$ 取所有非中心共轭类的代表，$C_G(x)=\{g\in G:gx=xg\}$ 是元素 $x$ 的中心化子。证明的关键是：要么中心中能取出一个阶为 $p$ 的正规子群，从而降到商群；要么类方程迫使某个较小的中心化子保留完整的 $p^a$ 因子。

\subsection{情形一：$p\mid |Z(G)|$}


中心 $Z(G)=\{x\in G:gx=xg\text{ 对所有 }g\in G\}$ 是有限阿贝尔群。由\textbf{阿贝尔群版本的 Cauchy 定理}（若素数 $p$ 整除有限阿贝尔群 $A$ 的阶，则 $A$ 中存在阶为 $p$ 的元素），$Z(G)$ 中存在阶为 $p$ 的元素 $z$。令

\[
N=\langle z\rangle=\{e,z,z^2,\dots,z^{p-1}\},\qquad |N|=p,
\]

其中 $\langle z\rangle$ 表示由 $z$ 生成的循环子群。

因为 $N\le Z(G)$，所以 $N\trianglelefteq G$，且

\[
|G/N|=\frac{|G|}{|N|}=p^{a-1}m.
\]

商群比 $G$ 小，归纳假设给出 $G/N$ 的一个 Sylow $p$-子群 $\overline P$，满足 $|\overline P|=p^{a-1}$。令自然映射 $\pi:G\to G/N$，取其原像

\[
P=\pi^{-1}(\overline P).
\]

由同态基本计数公式（或原像的阶公式）

\[
|P|=|N|\,|\overline P|=p\cdot p^{a-1}=p^a.
\]

于是 $P$ 就是 $G$ 的 Sylow $p$-子群。

\textbf{为什么这里要取原像？} 归纳只在商群 $G/N$ 中给出了一个子群；原像把它“抬回”到 $G$，并自动包含核 $N$，所以阶数多出恰好一个因子 $p$。

\subsection{情形二：$p\nmid |Z(G)|$}


（本节反复出现 $[G:H]=|G|/|H|$，即子群 $H\le G$ 的\textbf{指数}；$p\nmid n$ 读作“$p$ 不整除 $n$”。）

此时使用类方程。若对每个非中心代表 $x_i$ 都有

\[
p\mid [G:C_G(x_i)],
\]

那么右端的每一项都被 $p$ 整除，因而 $p\mid |Z(G)|$；这与当前情形矛盾。因此存在某个非中心元素 $x$，使得

\[
p\nmid [G:C_G(x)].
\]

记 $H=C_G(x)$。由于 $x\notin Z(G)$，$H$ 是 $G$ 的真子群，所以可以对 $H$ 使用归纳假设。又因为

\[
|G|=|H|[G:H],\qquad p\nmid [G:H],
\]

而 $|G|$ 的 $p$-部分是 $p^a$，可知 $|H|$ 仍含有完整的 $p^a$ 因子：存在 $k$ 使

\[
|H|=p^a k,\qquad p\nmid k.
\]

归纳假设给出 $H$ 中一个阶为 $p^a$ 的子群 $P$。把 $P$ 看成 $G$ 的子群，它的阶已经等于 $|G|$ 的最大 $p$-幂，因此 $P$ 是 $G$ 的 Sylow $p$-子群。

两种情形穷尽了所有可能性，存在性定理得证。

\subsection{课堂证明中容易错的地方}


\begin{itemize}
\item 归纳的第二种情形中，$H=C_G(x)$ 的阶是 $p^a k$，不是只剩下 $p^{a-1}$；因为选择的是一个不被 $p$ 整除的共轭类大小。
\item $C_G(x)$ 是\textbf{中心化子}，由“与元素 $x$ 交换”定义；后面对 Sylow 子群集合做共轭作用时出现的是\textbf{正规化子} $N_G(P)$，两者不能混用。
\item $G/N$ 的阶是 $p^{a-1}m$，不要把其中的 $m$ 漏掉。归纳得到的是商群中的阶 $p^{a-1}$ 子群，再取原像才得到 $G$ 中的阶 $p^a$ 子群。
\end{itemize}


\subsection{例：$S_4$ 中的 Sylow 子群阶数}


\[
|S_4|=24=2^3\cdot3.
\]

Sylow $2$-子群的阶必须是 $8$，Sylow $3$-子群的阶必须是 $3$。这里 $S_4$ 是 $\{1,2,3,4\}$ 上的对称群，轮换记号 $(1\,2\,3)$ 表示 $1\mapsto2\mapsto3\mapsto1$。例如

\[
\langle(1\,2\,3)\rangle=\{e,(1\,2\,3),(1\,3\,2)\}
\]

是一个 Sylow $3$-子群。存在性定理保证阶为 $8$ 的 Sylow $2$-子群也一定存在；至于这样的子群有多少个、彼此是否共轭，要等第二、第三定理处理。

\subsection{课堂补充：$S_n$ 的 Sylow $p$-子群怎么构造}


存在性定理说 Sylow 子群一定存在，但没说怎么找。在对称群里有一个可以手写的答案：\textbf{让置换只在自己的 $p$ 元块内部活动，再让它在块之间整体轮换}。课堂讲的就是这个构造，下面把它写成可复用的形式。

先做最简单的情形 $n=p^k$（$S_{p^k}$ 是 $p^k$ 个字母上的对称群），此时 $|S_{p^k}|$ 中 $p$ 的幂次是

\[
p^{1+p+p^2+\cdots+p^{k-1}}.
\]

把 $p^k$ 个点分成 $p$ 个块 $B_1,\ldots,B_p$，每块 $p^{k-1}$ 个点。递归地定义 $W_k\le S_{p^k}$：

\begin{itemize}
\item $k=1$：$W_1=\langle(1\,2\,\cdots\,p)\rangle$ 是 $p$ 阶循环群；
\item 从 $W_{k-1}$ 到 $W_k$：把 $W_{k-1}$ 的\textbf{每个生成元在每块内部同时作用}，再加上一个\textbf{把 $p$ 个块整体轮换}的置换。
\end{itemize}


这样得到的 $W_k$ 的阶是 $p^{k-1}\cdot|W_{k-1}|$，逐层展开正好是 $p^{1+p+\cdots+p^{k-1}}$，也就是 $S_{p^k}$ 的 Sylow $p$-子群。

\textbf{$S_4$，$p=2$。} $k=2$，$|S_4|$ 的 $2$-部分是 $2^3=8$。取两块 $\{1,2\}$、$\{3,4\}$：


\begin{center}
\begin{tabular}{ll}
生成元 & 含义\\ \hline
$(1\,2)$ & 第 1 块内部交换\\
$(3\,4)$ & 第 2 块内部交换\\
$(1\,3)(2\,4)$ & 两块整体轮换\\
$(1\,4)(2\,3)$ & 两块整体反向轮换
\end{tabular}
\end{center}


它们生成的子群阶为 $8$（直接乘出 $8$ 个元素即可验证），正是 $S_4$ 的 Sylow $2$-子群。注意这里 $4$ 个生成元比需要的多：$(1\,4)(2\,3)$ 可以由前三个乘出来，列出来只是为了看清“块内 $\times$ 块间”的结构。

\textbf{$S_6$，$p=3$（一个容易踩的坑）。} $6=3\cdot2$，$|S_6|$ 的 $3$-部分是 $3^2=9$，不是 $3\cdot3\cdot2=18$。如果天真地把“每块内部的 $3$-轮换”和“交换两块”都当成独立生成元，会得到一个 $18$ 阶子群，\textbf{它不是 $S_6$ 的 Sylow $3$-子群}（阶太大了）。正确的做法是让两块内部的 $3$-轮换相互关联，使整个子群只有 $9$ 个元素。这说明构造 Sylow 子群时，“阶必须整除 $|G|$ 的最大 $p$-幂”是必须回头检查的约束。

\subsection{参考对照}


\begin{itemize}
\item [DF] §4.5，Theorem 18（Sylow 定理及其归纳证明）；$S_n$ 的 Sylow 子群可用 $p$ 元块上的迭代\textbf{圈积}（wreath product，记作 $\wr$）描述——本节写出的递归构造就是它的具体展开，不需要预先知道圈积的定义。
\item [DS] 群作用与轨道计数的相关小节。该书正文没有 Sylow 定理，本节沿用课堂的类方程路线，并补全商群原像的阶数计算；$S_n$ 的构造按课堂补充整理。
\end{itemize}
